% Pillar-7 conclusion
\section{Conclusion}
\label{sec:p7-conclusion}
The companion Pillar-2 paper ended by refusing to over-read a diagnostic;
this audit explains what would be required to earn the control authority it
withheld. A one-line binomial model,
$S = p(1-p)(1 - h_G(p))$, predicts the ZVF phenomena Pillar~2 catalogued:
its exact expected value (T1, matched at the endpoints and interior of the
binned reward tensors), its monotone decline in group size (T2, confirmed
on two independent grids), and its coupling to gradient mass through
$p(1-p)$ (T3, directionally confirmed in an open backward pass at
$+0.71$). The same model exposes the diagnostic's structural limit---the
symmetry that aliases mastery with incapacity, visible as a U-shape across
368 runs---and supplies the repair: the pairwise-contrast density
$\mathrm{PCD} = \frac{G-1}{G}\E[p(1-p)]$, which survives the micro-jitter
that silently zeroes ZVF in any pipeline with a dense sub-reward.

The implemented callback and four-arm pilot establish feasibility, not
superiority. Adaptive $G$ reached the same held-out delta as Dr.~GRPO while
spending more rollouts, and the larger retrospective trace reveals that 92.3\%
of symmetric escalation fires occur on all-correct groups. Cross-prompt pooling
cannot stand in for a larger GRPO group, and a frozen-$p$ simulation cannot stand
in for closed-loop learning. The durable output is therefore a falsifiable design
rule---separate failed starvation from solved saturation, price interventions in
rollouts, and report the stack---plus a decisive next experiment: seed-paired,
fixed-token adaptive control against static $G{=}16$ and naive boundary
heuristics. Until that bakeoff succeeds, adaptive $G$ remains a proposal.
